% \begin{figure*}[ht]
  
%   \begin{minipage}{.46\linewidth}
%     \begin{center}
%       \Qcircuit @C=0.5em @R=0.2em @!R { \\
          \push{{x}_{0} }   & & \multigate{3}{x_k \cdot U}  & \qw   & \qw & &         & & \qw & \qw               & \qw                       & \qw      & \qw                       & \multigate{3}{U}              & \qw               & \qw      & \qw                         & \qw               & \qw & \\
          \push{{x}_{1} }   & & \ghost{x_k \cdot f}         & \qw   & \qw & &         & & \qw & \qw               & \qw                       & \qw      & \qw                       & \ghost{U}                      & \qw               & \qw      & \qw                         & \qw               & \qw & \\
          \push{\vdots}     & &                             &       &     & &\push{=} & &     &                   & \push{\vdots}             &          &  \push{\vdots}            &                                & \push{\vdots}     &          &                             & \push{\vdots}     &     & \\
          \push{{x}_{n-1} } & & \ghost{x_k \cdot f}         & \qw   & \qw & &         & & \qw & \qw               & \qw                       & \ctrl{1} & \qw                       & \ghost{U}                      & \qw               & \ctrl{1} & \qw                         & \qw               & \qw & \\
          \push{{x}_{n} }   & & \push{\rtoftarg}\qw\qwx[-1] & \qw   & \qw & &         & & \qw & \gate{\mathrm{H}} & \gate{\mathrm{T}}         & \targ    & \gate{\mathrm{T^\dagger}} & \targ\qwx[-1]    & \gate{\mathrm{T}} & \targ    & \gate{\mathrm{T^{\dagger}}} & \gate{\mathrm{H}} & \qw & \qw\\
}
                            

%     \end{center}
%     \caption{U-controlled RTOF gate}
%     \label{fig-urtof}
%   \end{minipage}
%   %\vspace*{0.25cm}
% \end{figure*}
% \vspace*{0.1cm}
\section{Preliminary}
\label{Sec:prelim}
\subsection{Quantum Boolean Circuits}

In this paper, we consider {\it quantum Boolean circuits}~\cite{bib-ddmf} that implement Boolean functions. We will be implementing them using the fault-tolerant set of Clifford gates, T gates, and S gates. 

We denote the primary inputs of a circuit by $x_1, x_2, \cdots,x_{n}$  when we have $n$ primary inputs. Each qubit state in a circuit can be represented by a Boolean function with respect to $x_1, x_2, \cdots,x_{n}$.

An example of such a quantum Boolean operation is the Toffoli gate depicted in
Fig ~\ref{fig-toffoli}. It inverts the basis state of the target bit (i.e.,
$\ket{1}$ to $\ket{0}$ and $\ket{0}$ to $\ket{1}$) when the states of the both
control bits are $\ket{11}$. This provides us a a way to realize the AND
operation in a quantum circuit.

\subsection{Relative-Phase Boolean Functions}

A typical decomposition of the Toffoli gate into the Clifford + T basis
consists of seven T gates as shown in Fig.~\ref{fig-toffoli}.
A Relative-Phase Toffoli (RTOF) gate~\cite{bib-rtof-maslov}, in contrast,
implements the same function using only four T gates as shown in
Fig.~\ref{fig-rtof}. However, it does this at the cost of a {\it relative phase},
which is an input dependent phase on the output.
Note the double circle on the target of Fig.~\ref{fig-rtof}: we will be
using this notation to denote a relative phase version of a logic
gate in this paper. One may notice that the version of the RTOF gate
that is used here is similar to the $n$-controlled-iX gate
from~\cite{bib-selinger-rtof} for $n=2$, with one less CNOT gate. As
noted in~\cite{bib-rtof-maslov}, removing this gate only affects the
phase, so it only transforms an RTOF gate into another RTOF gate with
the same T-count. We will be considering analogs of both types of
relative phase gatesin Section ~\ref{Sec:proposed}

As depicted in Fig.~\ref{fig-rtof}, an RTOF gate has three inputs; the input quantum basis states are $\ket{000}$
to $\ket{111}$. Similar to Toffoli gates, an RTOF gate swaps  $\ket{110}$ and $\ket{111}$. In addition to this
logic operation, RTOF gates add phases to some output states that the equivalent Toffoli gate does not. The added
phases are as shown in Fig.~\ref{table-rtof}. For example, when the input state is $\ket{110}$, a phase $+\pi/2$ is
added.

One of the results from~\cite{bib-rtof-maslov} is that when larger functions are broken down into compute/uncompute pairs, such as in the Barenco style Toffoli decompositions, it is almost always more advantageous to implement each ``half'' of the compute/uncompute pair into a relative phase construction. We will next go into one such construction.
%\vspace{-0.25cm}
\subsection{Boolean Logic Network}

A {\it Boolean logic network} is a simple Directed Acyclic Graph (DAG) whose vertices are primary inputs, primary outputs, and nodes which represent logic gates, whose arcs connect gates to inputs, outputs, and other gates.

A node has zero or more {\it fanins}, or incoming edges, and zero or more {\it fanouts} i.e. outgoing edges. The {\it primary inputs (PIs)} are nodes without fanins in the current network. The {\it primary outputs} are a subset of nodes of the network without fanouts. 

A logic function can be decomposed into such a Boolean logic network, and its nodes subsequently implemented as LUTs or logic gates. We see an example in Fig.~\ref{fig-lut-network}.

\begin{figure}[tbp]
    \centering
    \usetikzlibrary{positioning}
\definecolor{offwhite}{HTML}{F2EDED}
%% \tikzset{
%%     > = stealth,
%%     every node/.append style = {
%%         draw = none,
%%         text = offwhite
%%     },
%%     every path/.append style = {
%%         arrows = ->,
%%         draw = offwhite,
%%         fill = offwhite
%%     },
%%     hidden/.style = {
%%         draw = offwhite,
%%         shape = circle,
%%         inner sep = 1pt
%%     }
%% }
%% \tikz{
%%     \node (x) {$X$};
%%     \node[hidden] (z) [above right = of x] {$Z$};
%%     \node (y) [below right = of z] {$Y$};
%%     \path (x) edge (y);
%%     \path (z) edge (x);
%%     \path (z) edge (y);
%% }
\tikzset{
  > = stealth,
  every node/.append style = {
    draw      = black,
    shape     = circle,
    inner sep = 1
  },
  every path/.append style = {
        arrows = ->,
        draw = black,
        fill = offwhite
  },  
  input/.style = {
    draw = none
  }
}

\tikz{
  \node[input] (x1) at (0,0) {$x1$};
  \node[input] (x2) at (1,0) {$x2$};
  \node[input] (x3) at (2,0) {$x3$};
  \node[input] (x4) at (3,0) {$x4$};
  \node[input] (x5) at (4,0) {$x5$};
  \node[input] (y1) at (1,3) {$y1$};
  \node[input] (y2) at (3,4) {$y2$};
  \node (1) at (1,1) {$1$};
  \node (2) at (3,1) {$2$};
  \node (3) at (1,2) {$3$};
  \node (4) at (3,2) {$4$};
  \node (5) at (3,3) {$5$};      
  \path (x1) edge (3);
  \path (x2) edge (1);
  \path (x3) edge (1);
  \path (x4) edge (1);
  \path (x4) edge (2);
  \path (x5) edge (2);
  \path (x5) edge (5);
  \path (1)  edge (3);
  \path (1)  edge (4);    
  \path (2)  edge (4);
  \path (4)  edge (5);
  \path (3)  edge (y1);
  \path (5)  edge (y2);  
}

    \caption{A Boolean logic network}
    \label{fig-lut-network}
\end{figure}


\section{Existing Research and Motivation}
 \label{Sec:Exist}

\subsection{Using RTOF Gates to Decompose MCT Gates}
Existing research~\cite{bib-rtof-maslov}~\cite{bib-selinger-rtof} finds an advantage to using RTOF gates in synthesizing and decomposing logic functions, particularly Multiply-Controlled Toffoli Gates (MCTs).

In particular, Maslov proposes the decomposition depicted in Fig.~\ref{fig-maslov-decomp}~\cite{bib-rtof-maslov}. Here $U$ is a unitary operation implemented up to a relative phase.

Maslov further posits a way to construct $n$-input RTOF gates by replacing the middle CNOT gate in Fig.~\ref{fig-rtof} with an $n$-input RTOF gate, and then repeating recursively. However, as Amy et. al.~\cite{bib-amy-phase-state} note, the form of this RTOF gate generates a target-dependent phase which, when combined with the Hadamard gate at the end, adds the terms of its phase function to the Boolean function to the output, such that it takes the form $F = x_k \cdot F_{k-1} + F_{k-2}$.


\begin{figure}[tbp]
    \centering
    \scalebox{1}{\Qcircuit @C=0.5em @R=0.2em @!R { \\
  \push{\ket{{x}_{0}} } & & \multigate{4}{F} & \qw & &   & & \multigate{2}{F'} & \qw & \qw           & \qw & \multigate{2}{{F'}^{\dagger}} & \qw & \rstick{\ket{{x}_{0}}} \\
  \push{\ket{{x}_{1}} } & & \ghost{F}        & \qw & &   & & \ghost{F'}        & \qw & \qw           & \qw & \ghost{{F'}^{\dagger}}        & \qw & \rstick{\ket{{x}_{1}}}  \\
  \push{\ket{{x}_{2}} } & & \ghost{F}        & \qw & &   & & \ghost{F'}\qwx[1] & \qw &               & \qw & \ghost{{F'}^{\dagger}}\qwx[1] & \qw & \rstick{\ket{{x}_{2}}}  \\
  \push{\vdots}         & &                  &     & & = & & \push{\vdots}     &     & \push{\vdots} &     & \push{\vdots}                 &     & \rstick{\vdots }    \\       
  \push{\ket{{x}_{k}} } & & \ghost{F}        & \qw & &   & &  \qw              & \qw & \ctrl{2}      & \qw & \qw                           & \qw & \rstick{\ket{{x}_{k} }} \\
  \push{\ket{0} }       & & \qw              & \qw & &   & & \targ\qwx[-2]     & \qw & \control\qw   & \qw & \targ\qwx[-2]                 & \qw & \rstick{\ket{0} }  \\    
  \push{\ket{0} }       & & \targ\qwx[-2]    & \qw & &   & & \qw               & \qw & \targ         & \qw & \qw                           & \qw & \rstick{\ket{F} } \\
}

}
    \caption{Toffoli gate decomposition from~\cite{bib-rtof-maslov}}
    \label{fig-maslov-decomp}
\end{figure}


\begin{figure}[tbp]
    \centering
    \scalebox{1}{\Qcircuit @C=0.5em @R=0.2em @!R { \\
  \lstick{\ket{{x}_{1}}} & \qw  & \qw                 & \qw                 & \multigate{5}{U_3}      & \qw                & \qw                   & \qw                          & \qw                          & \qw                           & \ghost{x_1}  & \lstick{\ket{{x}_{1}}} \\
  \lstick{\ket{{x}_{2}}} & \qw  & \multigate{1}{U_1}  & \qw                 & \push{\dashedwire}\qw   & \qw                & \qw                   & \qw                          & \qw                          & \multigate{1}{U_1^{\dagger}}  & \ghost{x_2}  & \lstick{\ket{{x}_{2}}} \\
  \lstick{\ket{{x}_{3}}} & \qw  & \ghost{1}           & \qw                 & \push{\dashedwire}\qw   & \qw                & \qw                   & \qw                          & \qw                          & \ghost{U_1^{\dagger}}         & \ghost{x_3}  & \lstick{\ket{{x}_{3}}} \\
  \lstick{\ket{{x}_{4}}} & \qw  & \qw                 & \multigate{1}{U_2}  & \push{\dashedwire}\qw   & \qw                & \qw                   & \qw                          & \multigate{1}{U_2^{\dagger}} & \qw                           & \ghost{x_4}  & \lstick{\ket{{x}_{4}}} \\
  \lstick{\ket{{x}_{5}}} & \qw  & \qw                 & \ghost{2}           & \push{\dashedwire}\qw   & \qw                & \multigate{4}{U_5}    & \qw                          & \ghost{U_2^{\dagger}}        & \qw                           & \ghost{x_5}  & \lstick{\ket{{x}_{5}}} \\
  \lstick{\ket{0}     }  & \qw  & \targ\qwx[-3]\qw    & \qw                 & \ghost{3}\qw            & \multigate{1}{U_4} & \push{\dashedwire}\qw & \multigate{1}{U_4^{\dagger}} & \qw                          & \targ\qwx[-3]\qw              & \ghost{0}    & \lstick{\ket{0}     }    \\
  \lstick{\ket{0}     }  & \qw  & \qw                 & \targ\qwx[-2]\qw    & \qw                     & \ghost{4}          & \push{\dashedwire}\qw & \ghost{U_4^{\dagger}}        & \targ\qwx[-2]\qw             & \qw                           &  \qw            & \lstick{\ket{0}     }  \\
  \lstick{\ket{0}     }  & \qw  & \qw                 & \qw                 & \targ\qwx[-2]\qw        & \qw                & \push{\dashedwire}\qw & \qw                          & \qw                          & \qw                           & \ghost{y_1}  & \lstick{\ket{y_1}   }  \\      
  \lstick{\ket{0}     }  & \qw  & \qw                 & \qw                 & \qw                     & \targ\qwx[-2]\qw   & \ghost{5}             & \targ\qwx[-2]\qw             & \qw                          & \qw                           & \ghost{0}    & \lstick{\ket{0}     }    \\
  \lstick{\ket{0}     }  & \qw  & \qw                 & \qw                 & \qw                     & \qw                & \targ\qwx[-1]\qw      & \qw                          & \qw                          & \qw                           & \ghost{y_2}  & \lstick{\ket{y_2}     }  \\
}
}
    \caption{The quantum circuit implementing the Boolean logic network in Fig.~\ref{fig-lut-network} using the LHRS. }
    \label{fig-lut-qc}
\end{figure}

To summarize:
\begin{itemize}
\item Replacing pairs of computing and uncomputing gates with RTOF gate and its inverse will lead to T gate savings in most cases.
\item Using the above, it is possible to use clean ancilla to decompose an MCT with fewer T gates than the known methods.
\item By utilizing a relative UCNOT implementation, we get a recursive procedure to generate MCT gates up to a relative phase.
\end{itemize}



\subsection{LUT-based Hierarchical Reversible Logic Synthesis (LHRS)}
\label{Sec:lhrs}

Previous research~\cite{bib-lhrs-soeken} has detailed a way to leverage LUT-based synthesis of Boolean functions to create reversible logic networks that will be essential to quantum Boolean circuits. The overall flow maps each node in the Boolean network to a function-controlled NOT gate (UCNOT), a NOT gate controlled not just by the value of a variable but that of a Boolean function, and then cascades them together to create the output. 

The result of the method is shown in Fig~\ref{fig-lut-qc}. To get from Fig.~\ref{fig-lut-network} to Fig~\ref{fig-lut-qc} first, we take as input a LUT network synthesized from a Boolean expression using existing LUT-Network decomposition methods. Then this LUT network is parsed by the algorithm to create a quantum circuit of UCNOT gates. Any intermediate node is decomposed as two UCNOTs: one to compute, and the other to uncompute, acting on the same dirty ancilla. Any node that acts only on output qubits (output node i.e., an endpoint to the cascading logic) is given only one UCNOT.


Because the UCNOT gates act on ancilla, there is a need to uncompute them, so the flow next produces uncompute logic. This produces the compute/uncompute pair that a relative-phase based construction of the function can take advantage of. 

\begin{example}
  \label{ex-lut}
  We start at node 1. Since this is an intermediate node, we place a
  UCNOT on an ancilla, and place its inverse on the same ancilla to
  uncompute. We do the same for 2. For node 4, we cascade the ancilla
  from nodes 1 and 2 into it as input, and place its LUT, as well as
  its uncompute logic, on another ancilla. Finally, nodes 3 and 5 are
  output nodes, so we only generate their compute logic. 

\end{example}


%%% Local Variables:
%%% mode: japanese-latex
%%% TeX-master: "main"
%%% End:
